\vfill\pagebreak
\section{References}
\label{sec:mutability:references}

\token{discount} changes the elements of its caller's slice, because the
elements are shared. It cannot change a variable of its caller: the parameter
is a copy of the argument, and nothing it does reaches the variable the argument
came from. A function that exchanges the values of two integers needs more: it
must write into the variables of its caller.

A \textit{reference} is a parameter that holds the address of a variable of the
caller rather than a copy of its value. Reading the parameter reads that
variable, and assigning the parameter writes into it. It is declared with
\token{ref} before the name, and \token{ref mut} when the function may write.

\lstinputlisting[style=coloredverbatim, caption={Exchanging two variables, \textit{examples/mutability/swap.yr}}, label=lst:mutability:swap]{examples/mutability/swap.yr}
\vspace{-10pt}%
\begin{lstlisting}[style=bashVerb]
x = 2, y = 1
\end{lstlisting}

The parameters \token{a} and \token{b} of line~3 are references to integers
that the function may change. In the body, they are used as ordinary variables:
line~4 reads the value of the variable that \token{a} refers to, and line~5
writes into it. The call, on line~12, passes \token{ref x} and \token{ref y}:
the addresses of the two variables of \token{main}, which must be declared
\token{mut} for the function to change them.
Figure~\ref{fig:mutability:ref_parameter} shows the call.

\begin{figure}[h]
  \centering
  \begin{adjustbox}{max width=\linewidth}
    \begin{tikzpicture}
      \node[memhit] (x) at (0, 0) {2};
      \node[cflab, left=3mm of x] (xn) {\token{x}};
      \node[memcell] (y) at (0, -1.0) {2};
      \node[cflab, left=3mm of y] {\token{y}};
      \node[memframe, fit=(xn)(x)(y)] (main) {};
      \node[memtitle] at ([xshift=2mm]main.north west) {main};

      \node[memaddr] (a) at (0, -2.6) {};
      \node[cflab, left=3mm of a] (an) {\token{a}};
      \node[memaddr] (b) at (0, -3.6) {};
      \node[cflab, left=3mm of b] {\token{b}};
      \node[memcell] (old) at (0, -4.6) {1};
      \node[cflab, left=3mm of old] (on) {\token{old}};
      \node[memframe, fit=(an)(a)(old)] (swap) {};
      \node[memtitle] at ([xshift=2mm]swap.north west) {swap};

      \draw[mempoint] (a.center) to[out=0, in=0, looseness=1.2] (x.east);
      \draw[mempoint] (b.center) to[out=0, in=0, looseness=1.2] (y.east);

      \node[memregion] at ($(main.north) + (0, 5mm)$) {stack};
    \end{tikzpicture}
  \end{adjustbox}
  \caption{\label{fig:mutability:ref_parameter} The call
    \token{swap(ref x, ref y)} of Listing~\ref{lst:mutability:swap}, after
    \token{a = b;}. The parameters \token{a} and \token{b} hold the addresses of
    the variables of \token{main}, and \token{a = b;} wrote into \token{x}.
    \token{old} still holds the first value of \token{x}, which
    \token{b = old;} writes into \token{y} next.}
\end{figure}


As \token{alias} does, the \token{ref} of the call tells a reader of
\token{main} that the call may change the variables it names. It is required.
\token{swap(x, y)} is refused with \errmsg{E4226}{the call operator is not
  defined for \errhole{} and \errhole{}}, the note under it saying
\textit{cannot pass value of type mut i32 as ref}. A reference refers to a
variable, so its argument must be one: \token{swap(ref x, ref 3)} is refused
with \errmsg{E4151}{not a lvalue}, and a variable declared without \token{mut}
cannot be passed to \token{ref mut}, since the function would change a
constant.

\subsection{Giving the caller's variable a new slice}

A reference can refer to a variable of any type, including a slice. Through
it, a function can give its caller's variable a new slice, which a
\token{dmut} parameter cannot do (Section~\ref{sec:mutability:parameters}).

\lstinputlisting[style=coloredverbatim, caption={Appending to the caller's slice, \textit{examples/mutability/history.yr}}, label=lst:mutability:history]{examples/mutability/history.yr}
\vspace{-10pt}%
\begin{lstlisting}[style=bashVerb]
[18, 21, 19]
\end{lstlisting}

\tokennolst{history \mtildeascii{}= [value]}, on line~4, gives the variable a
new slice, one element longer. With a parameter declared \token{mut history:
  [i32]}, this would have changed a copy, and \token{main} would have printed an
empty slice; the compiler refuses such a parameter anyway, with the
\errmsg{E4138}{a parameter cannot be mutable, if it is not a reference or does
  not borrow mutable data} of Section~\ref{sec:mutability:parameters}. The
variable of \token{main} is \token{mut}: only the variable changes, not the
elements already in it.

\subsection{Constant references}

A reference without \token{mut} gives read access to the caller's variable. It
is useful for a large value stored in its variable, such as an array of a
thousand elements: a parameter of type \token{[i32 ; 1000]} copies the thousand
elements into the frame at each call, a reference copies an address.

\begin{lstlisting}[style=coloredverbatim]
fn last(ref values: [i32 ; 1000])-> i32 {
  values[999]
}

fn main() {
  let big = [7 ; 1000];
  println(last(big));
}
\end{lstlisting}

The call writes nothing: the function cannot change the variable, so there is
nothing for the reader to be warned of. The \token{ref} keyword in a call is
written only for a reference that may change, and is refused elsewhere, as
\token{alias} is. A slice needs no constant reference: it is already two
numbers, and a plain parameter shares its elements.

\subsection{References to elements}

The argument of a reference can be an element of an array or of a slice, which
is a place in memory as a variable is. \token{swap(ref t[0], ref t[2])}
exchanges the first and the third elements of \token{t}, a slice declared
\token{dmut}. To change the elements of an array of its caller, a function
takes a reference to the whole array, declared \token{ref dmut}: \token{ref mut}
alone lets it give the array a new value, but not change its elements, for the
reason of Section~\ref{sec:mutability:levels}.

\begin{lstlisting}[style=coloredverbatim]
fn clear(ref dmut counts: [i32 ; 6]) {
  for i in 0 .. 6 {
    counts[i] = 0;
  }
}

fn main() {
  let dmut counts = [1, 2, 3, 4, 5, 6];
  clear(ref counts);
  println(counts);
}
\end{lstlisting}

\subsection{Choosing a parameter}

\begin{center}
  \begin{adjustbox}{max width=\linewidth}
    \begin{tabular}{|l|l|l|}
      \hline
      The function & Parameter & Call\\[0pt]
      \hline
      \hline
      reads its argument & \token{x: T} & \token{f(x)}\\[0pt]
      reads a large array without copying it & \token{ref x: [T ; n]} & \token{f(x)}\\[0pt]
      changes the elements of a slice or a map & \token{dmut x: [T]} & \token{f(alias x)}\\[0pt]
      gives the caller's variable a new value & \token{ref mut x: T} & \token{f(ref x)}\\[0pt]
      changes the elements of the caller's array & \token{ref dmut x: [T ; n]} & \token{f(ref x)}\\[0pt]
      \hline
    \end{tabular}
  \end{adjustbox}
\end{center}

A reference is only a parameter, or a local variable declared
\token{let ref}, which Part II describes (Section~\ref{sec:memory:ref_variable}).
It cannot be stored in another value, nor returned by a function: the variable
it refers to may be in the frame of the function, which disappears when the
function returns.
